\chapter{Kéo Lại Năng Lượng Cuối Buổi}
\chapterintro{Cuối buổi không có nghĩa là hết học}{Đây là lúc lớp dễ mỏi nhất, nhưng cũng là lúc vài nhịp đúng có thể tạo ấn tượng mạnh nhất.}

\section{Đứng lên trả lời ngắn}
\begin{tinhhuong}{Tình huống}
Cuối buổi, lớp ngồi đã lâu và năng lượng tụt rất nhanh.
\end{tinhhuong}
\begin{goiydungngay}
Mời 3 em đứng lên, mỗi em nói 1 câu ngắn về bài rồi ngồi xuống ngay.
\begin{itemize}
\item Đổi tư thế cơ thể giúp tỉnh lớp.
\item Chỉ dùng rất ngắn để không mất thêm sức.
\item Hợp cuối buổi sáng hoặc tiết 4, tiết 5.
\end{itemize}
\end{goiydungngay}
\begin{cauchot}
Muốn lớp tỉnh lại, đôi khi chỉ cần thay đổi tư thế một chút.
\end{cauchot}
\ngancach

\section{Một tiếng cười trước khi chốt bài}
\begin{tinhhuong}{Tình huống}
Lớp mệt thấy rõ, nhưng bài vẫn còn một đoạn cần kết lại cho gọn.
\end{tinhhuong}
\begin{goiydungngay}
Cho một câu rất nhẹ: \textit{``Ai đang chiến đấu với cơn buồn ngủ thì gật đầu nhẹ cho thầy cô biết.''}
\begin{itemize}
\item Thừa nhận thực tế giúp lớp dễ chịu.
\item Không kéo sang đùa quá lâu.
\item Cười một chút rồi chốt nhanh phần còn lại.
\end{itemize}
\end{goiydungngay}
\begin{cauchot}
Cuối buổi, sự cảm thông thường hiệu quả hơn quát tháo.
\end{cauchot}
\ngancach

\section{Chốt bằng việc gần nhất}
\begin{tinhhuong}{Tình huống}
Cuối buổi đầu óc học sinh không còn hợp với những câu hỏi quá rộng.
\end{tinhhuong}
\begin{goiydungngay}
Hỏi đúng một điều rất gần: \textit{``Nếu về nhà chỉ làm một việc liên quan đến bài này, em làm gì?''}
\begin{itemize}
\item Câu hỏi gần giúp lớp đáp nhanh.
\item Kết nối bài học với hành động cụ thể.
\item Chốt tiết gọn mà không rơi.
\end{itemize}
\end{goiydungngay}
\begin{cauchot}
Cuối buổi, câu hỏi càng gần thực tế thì càng giữ được lớp.
\end{cauchot}
\ngancach

\section{Đếm ngược một phút cuối có mục tiêu}
\begin{tinhhuong}{Tình huống}
Cuối buổi lớp rất dễ tan trước khi nội dung thật sự khép lại.
\end{tinhhuong}
\begin{goiydungngay}
Thông báo: \textit{``Mình còn đúng một phút để giữ lại điều quan trọng nhất.''}
\begin{itemize}
\item Việc đếm ngược làm lớp chú ý lại.
\item Trong một phút ấy chỉ làm đúng một việc.
\item Có thể là chốt ý, đặt hẹn, hoặc hỏi một câu cuối.
\end{itemize}
\end{goiydungngay}
\begin{cauchot}
Khi lớp biết thời gian còn ít nhưng rõ việc, các em dễ dồn lại năng lượng hơn.
\end{cauchot}
\ngancach

\section{Nối bài với việc tan trường}
\begin{tinhhuong}{Tình huống}
Cuối buổi, suy nghĩ của học sinh đã chạm tới cổng trường và chuyện về nhà.
\end{tinhhuong}
\begin{goiydungngay}
Hỏi: \textit{``Tan học xong, bài này còn dùng được ở việc nào gần em nhất?''}
\begin{itemize}
\item Câu hỏi gần đời sống nên dễ kéo phản hồi.
\item Học sinh thấy bài chưa bị bỏ lại trong lớp.
\item Rất hợp để chốt những tiết cuối ngày.
\end{itemize}
\end{goiydungngay}
\begin{cauchot}
Bài học bớt rơi khi nó đi cùng học sinh ra khỏi lớp.
\end{cauchot}
